\documentclass[ti,nofooter]{own-haw-page}
\usepackage[german]{babel}
\usepackage{color}
\usepackage{enumitem}
\usepackage{nameref}
\usepackage{url}
\usepackage{seqsplit}
\usepackage[utf8]{inputenc}
\usepackage[hidelinks]{hyperref}
\usepackage[many]{tcolorbox}
\usepackage{tabto}

\newcommand{\raus}[1]{}
\newcommand{\aufg}{\marginpar{{\em Protokoll!}}}
\newcommand{\back}{\ensuremath{\backslash}}
\newcommand{\pfeil}{\ensuremath{\rightarrow}}
\newcommand{\platzlinie}[1]{%
~\\[#1]
~\\
%$\underline{\hfill}$~\\
\hrule~\\
}

\definecolor{lightblue}{rgb}{0.8, 0.9 ,1}

\newcommand{\arrow}{{\textgreater{}\textgreater} }

% Tilde%%%%%
\newcommand{\addtildealt}{\textasciitilde{}}
% schicker, funktioniert noch nicht:
%\newcommand{\addtilde}{{\Large\textasciitilde{}}}

\newcommand{\addtildealts}{{\fontfamily{ptm}\selectfont\textasciitilde}{}}

\newcommand{\addtilde}{\raisebox{-0.3em}\textasciitilde{}}

% Abstand
\newcommand{\addspace}[1]{\vspace{#1}}

% Codezeile mit $ beginnend
\newcommand{\command}[1]{%
  \vspace{10pt}
  {\ttfamily \fcolorbox{black}{lightblue}{
      \parbox{\textwidth} {
        \vspace{-10pt}
        \def\labelitemi{\$}
        \begin{itemize}[leftmargin=3.5ex, itemsep=-1ex, rightmargin=4ex]
          #1
          \vspace{-10pt}
      \end{itemize}}
  }}\vspace{10pt}}

%%%%%%%%%%%%%%%%%%%%%%%%%%%%%%%%%%%%%%%%%%
% Neues Design (tcolorbox)
\tcbset{settings/.style={
    fontupper=\ttfamily,
    breakable,
    enhanced, 
    colback=lightblue, 
    arc=0mm, 
    boxrule=0.5pt,
    enlarge bottom by=0.5ex,
    enlarge top by=1ex,  
    left=1.5ex,
    right=0ex,
    top=0.4ex,
    bottom=0.4ex}}

% Nichtnummerierter Code (neu)
\newcommand{\codespace}[1]{
  \begin{tcolorbox}[settings]{}
    #1
  \end{tcolorbox}
  \vspace{10pt}
}

% Codezeile mit $ beginnend (neu, unbenutzt)
\newcommand{\codedollar}[1]{
  \begin{tcolorbox}[settings]{}
    \def\labelitemi{\$}
    \begin{itemize}[leftmargin=3ex, itemsep=-1ex]
      #1
    \end{itemize}
  \end{tcolorbox}
}

% Unbenutzt
\newcommand{\codenumber}[1]{
  \begin{tcolorbox}[settings]{}
    \begin{enumerate}[leftmargin=3ex, label={\arabic*}, itemsep=-1ex]
      #1
    \end{enumerate}
  \end{tcolorbox}
}
%%%%%%%%%%%%%%%%%%%%%%%%%%%%%%%%%%%%%%%%%%

% Nichtnummerierter Code (alt)
\newcommand{\codebox}[1]{%
  \vspace{10pt}
  {\ttfamily \fcolorbox{black}{lightblue}{
      \parbox{\textwidth} {
        \vspace{-10pt}
        \def\labelitemi{}
        \begin{itemize}[leftmargin=1ex, itemsep=-1ex, rightmargin=4ex]
          #1
          \vspace{-10pt}
      \end{itemize}}
  }}\vspace{10pt}}

% Normaler nummerierter Code
\newcommand{\code}[1]{
  \vspace{10pt}
  %{\setlength{\fboxrule}{1pt} % Rahmendicke definieren
  {\ttfamily \fcolorbox{black}{lightblue}{
      \parbox{\textwidth} {
        \vspace{-10pt}
        \begin{enumerate}[leftmargin=3.5ex, label={\arabic*}, itemsep=-1ex, rightmargin=4ex]
          #1
          \vspace{-10pt}
      \end{enumerate}}
  }}\vspace{10pt}}
%}

\sloppy 



\begin{document}
%\date{16.~April~2020}
%\date{24.~April~2020}
\date{10.~Dezember~2021}


%\setcounter{section}{0}
%\setcounter{section}{1}
\setcounter{section}{5}
\newcounter{dummymcpdsv} 
\setcounter{dummymcpdsv}{\value{section}}
\addtocounter{dummymcpdsv}{1}


\begin{page}

\begin{Large}
{\centering \bfseries \"{U}bungsblatt \arabic{dummymcpdsv}\\ zur Veranstaltung  
\glqq RobE -- Einf\"{u}hrung in die Robotik''\\}
\end{Large}
~\\[-0.5cm]


\section{Trajektorienplanung eines Manipulators}

In der \"{U}bung \arabic{dummymcpdsv} geht es um die Trajektorienplanung des Manipulators Panda der Firma Franka Emika mittels des MoveIt!-Framework in ROS. \\


\textbf{Zur Bearbeitung der \"{U}bungsaufgaben:} Bearbeiten Sie auf jeden Fall
\textit{alle} \"{U}bungsaufgaben. Ausgenommen hiervon sind lediglich die mit \glqq
optional'' gekennzeichneten Textstellen. Lesen Sie sich die Aufgaben gut durch.
Sollten Sie eine Aufgabe nicht l\"{o}sen k\"{o}nnen, so beschreiben Sie
zumindest, wie weit Sie gekommen sind und auf welche Weise Sie vorgegangen sind.
%
Aufgaben mit der Randbemerkung \textit{``Protokoll!}'' sind 
%Die Aufgaben sind direkt 
im Protokoll zu beantworten.
Als L\"{o}sung der Aufgaben ist %entweder
%alternativ
(je nach Aufgabe)
Programmcode abzugeben
(kommentiert) oder ein kurzer Text.


\subsection{Vorbereitungsaufgaben}

%\begin{enumerate}
%\item Schauen Sie sich einmal die Beschreibungen zu den ROS-packages 
%\end{enumerate}

-- keine --

\subsection{Praktikumsaufgaben}

\begin{enumerate}

\item Bearbeiten Sie die Schritte \textbf{\nameref{sec:setup}} (Anhang \ref{sec:setup}).

\item Bearbeiten Sie die Schritte \textbf{\nameref{sec:start}} (Anhang \ref{sec:start}). Welche verschiedenen Bereiche geh\"{o}ren zum \textit{MotionPlanning} und sind nun in \textit{RVIZ} zu sehen? \"{U}berlegen Sie, warum diese Aufteilung notwendig ist. \aufg\

\item W\"{a}hlen Sie nun eine beliebige Zielstellung (\textit{Goal state}) und planen Sie eine Trajektorie. W\"{a}hlen Sie nun den Algorithmus \textit{RRTstarkConfigDefault} und f\"{u}hren Sie eine weitere Planung aus. Wiederholen Sie das gleich mit \textit{BFMtkConfigDefault}. Welcher Einfluss auf den Planungsprozess ist zu beobachten? \aufg\

\item W\"{a}hlen Sie nun die Option \textit{Use Carthesian Path} aus und planen Sie eine neue Trajektorie Wie hat diese ver\"{a}ndert? Visualisieren Sie die geplante Trajektorie in einer Schrittgr\"{o}ße (\textit{step size}) von 10 und planen Sie die Trajektorie mit ein- und ausgeschalteter Option. Was f\"{a}llt nun auf? \aufg\

\item Bearbeiten Sie die Schritte \textbf{\nameref{sec:commander}} (Anhang \ref{sec:commander}). Bewegen Sie nun den Endeffektor 20 cm nach oben. Was zeichnet die Bewegung bez\"{u}glich der inversen Kinematik aus?\aufg\

\item Positionieren Sie den Endeffektor zwischen zwei Boxen in Bodenn\"{a}he und setzen Sie das n\"{a}chste Ziel zwischen die n\"{a}chsten Boxen. Planen Sie nun die Trajektorie. Welches Verhalten wird beobachtet?   \aufg\

\item Bearbeiten Sie die Schritte \textbf{\nameref{sec:pickplace}} (Anhang \ref{sec:pickplace}). Beschreiben Sie die verschiedenen Schritten des \textit{Pick-and-Place}-Prozesses. \aufg\

\item \"{O}ffnen Sie den Quellcode zu der \textit{Pick-and-Place}-Demonstration. Beschreiben Sie Hauptfunktionen des Prozesses und ordnen Sie diesen den zuvor beschriebenen Schritten zu. \aufg\

\item \"{A}ndern Sie nun die x-Koordinate des zu platzierenden Objekts in der \textit{place}-Funktion zu 0.2 und bauen Sie den \textit{catkin workspace} erneut. Was ist bei erneuter Ausf\"{u}hrung der Demonstration zu beobachten?  \aufg\

\end{enumerate}


\appendix

%%%%%%%%%%%%%%%%%%%%%%%%%%%%%%%%%%%%%%%%%%%%%%%%%%%%%
%%%%%%%%%%%%%%%%%%%%%%%%%%%%%%%%%%%%%%%%%%%%%%%%%%%%%
%%%%%%%%%%%%%%%%%%%%%%%%%%%%%%%%%%%%%%%%%%%%%%%%%%%%%

\section{Installation and setup of MoveIt!}
\label{sec:setup}

The detailed installation process can be found here: \\

\url{https://ros-planning.github.io/moveit_tutorials/doc/getting_started/getting_started.html} \\

Please download the zipped package \textit{moveit\_ws.zip} to your \textit{home} directory and unpack. It can be found in \textbf{Teams}. This will be your workspace folder. 

\code{
\item echo {\dq source /opt/ros/noetic/setup.bash\dq} >""> \addtilde/.bashrc
\item echo {\dq export PATH=/usr/lib/ccache:\$PATH\dq} >""> \addtilde/.bashrc
\item source \addtilde/.bashrc
}

Finally, build the \textit{source} folder with:

\code{
\setcounter{enumi}{3}
\item sudo apt-get update
\item sudo apt-get install ros-noetic-moveit \\ros-noetic-panda-moveit-config ros-noetic-moveit-visual-tools
\item mkdir -p moveit\_ws/src
\item cd moveit\_ws/src
\item git clone https://github.com/ros-planning/moveit\_tutorials
\item git clone https://github.com/ros-planning/moveit\_task\_constructor
\item rosdep update
\item rosdep install -{}-from-paths . -{}-ignore-src -{}-rosdistro noetic -y
\item catkin build
\item echo {\dq source \addtilde/moveit\_ws/devel/setup.bash\dq} >""> \addtilde/.bashrc
}


%%%%%%%%%%%%%%%%%%%%%%%%%%%%%%%%%%%%%%%%%%%%%%%%%%%%%
\section{First steps with MoveIt!}
\label{sec:start}

Let's start with taking a look into a simple demonstration in \textit{RVIZ}. \\

Launch the simple demo:

\code{
\setcounter{enumi}{13}
\item source ~/.bashrc
\item roslaunch panda\_moveit\_config demo.launch rviz\_tutorial:=true
}

To visualize the motion planning, add the \textit{MotionPlanning}-view via the button \textit{Add}. On the bottom left you will have a new window. Choose the tab \textit{Planning} and select \textit{panda\_arm} as \textit{planning group} in order to manually move the end-effectors target position. Press \textit{Plan} to start the motion planning. After that, select \textit{Plan and Execute}. There are plenty of path planning algorithms. You can find them under the tab \textit{Context}. \\

You can visualize every step of the planned path by going to \textit{Planned path} in  \textit{MotionPlanning} in \textit{Displays} and select the option \textit{Show Trail}. By modifying \textit{Trail Step Size} you can skip a specified number of steps.


%%%%%%%%%%%%%%%%%%%%%%%%%%%%%%%%%%%%%%%%%%%%%%%%%%%%%
\section{Planning trajectories the MoveIt!-commander}
\label{sec:commander}

\textit{MoveIt!} comes with an C++ and Python Interface that provides a set of functions for motion control. A Python-based \textit{command line tool} allows to access these function in a very simple way. \\

To begin, the MoveIt! Panda demonstration has to be launch:

\code{
\setcounter{enumi}{15}
\item roslaunch panda\_moveit\_config demo.launch
}

The \textit{MoveIt!-Commander} can be launch with:

\code{
\setcounter{enumi}{16}
\item rosrun moveit\_commander moveit\_commander\_cmdline.py
}

To view all the functionalities, type \textit{help}. \\

The \textit{commander} has to be connected to the right \textit{move group node}, in order to send motion commands. This can be done by typing:

\code{
\setcounter{enumi}{17}
\item use panda\_arm
}

You can now move the end-effector by sending commands.Try:

\code{
\setcounter{enumi}{18}
\item go <up, down, left, right> <distance in meters>
}

The current state can be check with

\code{
\setcounter{enumi}{19}
\item current
}

and saved to a variable

\code{
\setcounter{enumi}{20}
\item rec <state name>
}

A trajectory can then be planned and executed to that saved state in one step with

\code{
\setcounter{enumi}{21}
\item go <state name>
}

or two steps

\code{
\setcounter{enumi}{22}
\item plan <state name>
\item execute
}

Furthermore, a numbered list of commands can be saved to a file and then loaded within the \textit{commander}. \\

Let's make a command list. Go to:

\code{
\setcounter{enumi}{24}
\item cd \addtilde/moveit\_ws/src
\item gedit command\_list
}

and save this example:

\codebox{
\item use panda\_arm \\
go up 0.3 \\
rec state1 \\
go rotate 0 0 1 \\
go up 0.4 \\
rotate 1 0 0 \\
go state1 \\
exit
}

Save this file as \textit{command\_list} in your workspace source folder. To execute the command list, open the \textit{commander} and type:

\code{
\setcounter{enumi}{26}
\item load command\_list
}

We can load collision objects into the planning scene:

\code{
\setcounter{enumi}{27}
\item rosrun moveit\_tutorials collision\_scene\_example.py cluttered
}

These objects are taken into account when planning a path to avoid collision.


%%%%%%%%%%%%%%%%%%%%%%%%%%%%%%%%%%%%%%%%%%%%%%%%%%%%%
\section{Pick-and-place motion planning}
\label{sec:pickplace}

Now we are going to play through the \textit{Pick-and-Place}-demonstration and take a look into its process. \\

Begin as usual:

\code{
\setcounter{enumi}{28}
\item source \addtilde/moveit\_ws/devel/setup.bash
\item roslaunch panda\_moveit\_config demo.launch
}

Start the \textit{Pick-and-Place}-demonstration: \\

\code{
\setcounter{enumi}{29}
\item rosrun moveit\_tutorials pick\_place\_tutorial
}

The demonstration will automatically start. \\

The source code to this demonstration can be found under:

\code{
\setcounter{enumi}{30}
\item \addtilde/moveit\_ws/src/moveit\_tutorials/doc/pick\_place/src
\item gedit pick\_place\_tutorial.cpp
}

For example, the target position and pose of the placed object can be modified within the \textit{place}-function. Also, collision objects can be added or modified in the \textit{addCollisionObjects}-function. For that, the number of objects has to be changed and the object's name, dimension, pose and primitive have to be added. To grasp the newly added object, the function \textit{pick} and \textit{place} has to be adapted to it as wells. \\

After changes have been made, the \textit{catkin workspace} has to be rebuild, using:

\code{
\setcounter{enumi}{32}
\item cd \addtilde/moveit\_ws
\item catkin build
}


%%%%%%%%%%%%%%%%%%%%%%%%%%%%%%%%%%%%%%%%%%%%%%%%%%%%%
%%%%%%%%%%%%%%%%%%%%%%%%%%%%%%%%%%%%%%%%%%%%%%%%%%%%%
%%%%%%%%%%%%%%%%%%%%%%%%%%%%%%%%%%%%%%%%%%%%%%%%%%%%%


\end{page}

\end{document}
